\documentclass[11pt,a4paper]{article}

% ── Self-contained preamble (does not depend on main.tex's preamble.tex) ────
\usepackage[utf8]{inputenc}
\usepackage[T1]{fontenc}
\usepackage{lmodern}
\usepackage{amsmath, amssymb, amsthm}
\usepackage{graphicx}
\usepackage{booktabs}
\usepackage{longtable}
\usepackage{hyperref}
\usepackage{url}
\usepackage{geometry}
\usepackage{xcolor}
\usepackage{listings}
\usepackage{bm}

\geometry{margin=1in}

\hypersetup{
    colorlinks=true,
    linkcolor=blue,
    filecolor=magenta,
    urlcolor=cyan,
}

\theoremstyle{definition}
\newtheorem{definition}{Definition}[section]
\newtheorem{theorem}{Theorem}[section]

\newcommand{\dd}{\mathrm{d}}
\newcommand{\Ncal}{\mathcal{N}}
\newcommand{\Ecal}{\mathbb{E}}
\newcommand{\Pcal}{\mathbb{P}}
\newcommand{\Reals}{\mathbb{R}}

% Code listing style
\definecolor{codegreen}{rgb}{0,0.6,0}
\definecolor{codegray}{rgb}{0.5,0.5,0.5}
\definecolor{codepurple}{rgb}{0.58,0,0.82}
\definecolor{backcolour}{rgb}{0.97,0.97,0.95}
\lstdefinestyle{code}{
    backgroundcolor=\color{backcolour},
    commentstyle=\color{codegreen},
    keywordstyle=\color{magenta},
    numberstyle=\tiny\color{codegray},
    stringstyle=\color{codepurple},
    basicstyle=\ttfamily\footnotesize,
    breakatwhitespace=false,
    breaklines=true,
    captionpos=b,
    keepspaces=true,
    numbers=left,
    numbersep=5pt,
    showspaces=false,
    showstringspaces=false,
    showtabs=false,
    tabsize=2,
    language=Python,
}
\lstset{style=code}

\title{FSA-v5 Technical Implementation Guide \\
\large A Reference for Engineers Porting the Model into \texttt{smc2fc}}
\author{FSA Model Development Team}
\date{\today}

\begin{document}

\maketitle

\begin{abstract}
\noindent
This document is the definitive technical reference for the FSA-v5
physiological model as implemented in \path{models/fsa_high_res/}. Its
audience is the research engineer who must integrate the model into the
\texttt{smc2fc} inference and control framework. It covers the full
mathematical specification (latent SDE + observation likelihood), the
v5 deconditioning extension that produces the closed-island basin
topology, the chance-constrained probabilistic cost functional, the
numerical integration protocol, the parameter freeze policy and its
Fisher-information justification, end-to-end test scenarios on
synthetic data, and a list of bugs and gotchas surfaced during the v5
implementation pass. A complete parameter reference is included as an
appendix.

This document deliberately omits the historical narrative arcs of v2,
v3, and v4 that appear in the longer modelling notes (\path{main.tex}).
For background motivation, that document is the reference; for
implementation, this document is the reference.
\end{abstract}

\tableofcontents
\newpage

% =============================================================================
\section{Overview} \label{sec:overview}
% =============================================================================

\subsection{What this guide covers}

The FSA-v5 model is a 6-dimensional stochastic differential equation
(SDE) with five observation channels. Forward simulation, Bayesian
inference, and closed-loop control of this model are all supported by
the code in \path{models/fsa_high_res/}. This guide explains:

\begin{itemize}
  \item The mathematical specification --- every latent equation,
        every observation channel, every parameter.
  \item The closed-island basin topology of v5 --- the model's
        prediction that there are three physiologically meaningful
        steady states (mono-stable healthy, sedentary collapse,
        over-training collapse).
  \item The chance-constrained probabilistic cost functional that
        defines what ``optimal training'' means under v5.
  \item The numerical integration protocol (15-minute bin grid, no
        sub-stepping, state-clipping) and its identifiability
        consequences.
  \item Which parameters are estimated and which are frozen, and the
        Fisher-information argument behind that choice.
  \item Three constant-$\Phi$ deterministic test scenarios that an
        engineer can run during the port to verify the implementation
        works as expected.
  \item Common bugs and gotchas surfaced during the v5 implementation
        pass --- read this section before you start porting.
\end{itemize}

\subsection{What this guide deliberately omits}

This guide does not cover the historical evolution from v2 through v4,
the FSA-v3 bimodal extension, the v4 Variable-Dose extension, or the
gradient-OT control formulation. Those are in \path{main.tex} for
background. This document treats v5 as the canonical model.

\subsection{Files and the integration boundary} \label{sec:files}

The complete v5 implementation lives in
\path{models/fsa_high_res/}. The integration boundary --- everything
\texttt{smc2fc} should import from --- is the package's
\texttt{\_\_init\_\_.py}, which re-exports the following symbols:

\begin{center}\small
\begin{tabular}{lll}
\toprule
Symbol & Source file & Role \\
\midrule
\texttt{HIGH\_RES\_FSA\_V5\_MODEL}      & \path{simulation.py} & Forward-sim \texttt{SDEModel} \\
\texttt{HIGH\_RES\_FSA\_V5\_ESTIMATION} & \path{estimation.py} & Inference \texttt{EstimationModel} \\
\texttt{StepwisePlant}                  & \path{_plant.py}     & Closed-loop synthetic plant \\
\texttt{evaluate\_chance\_constrained\_cost} & \path{control_v5.py} & v5 cost (eqs.\ \ref{eq:chance-constraint}, \ref{eq:v5-min-effort}) \\
\texttt{find\_A\_sep\_v5}               & \path{control_v5.py} & Separatrix root-finder \\
\texttt{TRUTH\_PARAMS\_V5}              & \path{_dynamics.py}  & v5 truth parameter dict \\
\texttt{DEFAULT\_PARAMS\_V5}            & \path{simulation.py} & v5 truth + observation coefficients \\
\texttt{DEFAULT\_INIT}                  & \path{simulation.py} & 6D initial state defaults \\
\texttt{drift\_jax}                     & \path{_dynamics.py}  & Canonical drift (single source of truth) \\
\texttt{diffusion\_state\_dep}          & \path{_dynamics.py}  & State-dependent diffusion \\
\texttt{BINS\_PER\_DAY, DT\_BIN\_DAYS}  & \path{simulation.py} & Time-grid constants \\
\texttt{circadian, circadian\_jax}      & \path{simulation.py} & Cosine circadian regressor \\
\texttt{build\_control\_spec\_v5}       & \path{control.py}    & Gradient-OT spec (legacy/comparison) \\
\bottomrule
\end{tabular}
\end{center}

A v4 alias (\texttt{HIGH\_RES\_FSA\_V4\_*}) is kept for back-compat with
existing tests; it is numerically identical to v5 with the Hill
deconditioning amplitudes pinned at zero.

\subsection{How to read this guide}

Sections 2--3 are the model specification: read them top to bottom in
order. Sections 4--5 give the v5 contributions (basin geometry +
probabilistic cost) that distinguish v5 from prior versions. Sections
6--7 are operational --- how to integrate and what to estimate.
Sections 8--9 are practical --- test scenarios and bugs to anticipate.
The appendix is a parameter reference. Code snippets throughout cite
specific files in the package.


% =============================================================================
\section{Latent dynamics: the FSA-v5 SDE} \label{sec:latent}
% =============================================================================

\subsection{State vector}

The latent state at time $t$ is the 6-vector
\begin{equation}
y(t) \;=\; \big( B(t),\; S(t),\; F(t),\; A(t),\; K_{FB}(t),\; K_{FS}(t)\big)^T,
\end{equation}
with the components defined as follows.

\begin{center}\small
\begin{tabular}{llll}
\toprule
Symbol & Range & Physiological role & Diffusion family \\
\midrule
$B$    & $[0, 1]$        & Aerobic fitness (Banister chronic capacity)         & Jacobi, $\sqrt{B(1-B)}$ \\
$S$    & $[0, 1]$        & Strength capacity (Banister chronic capacity)        & Jacobi, $\sqrt{S(1-S)}$ \\
$F$    & $[0, \infty)$   & Unified fatigue pool                                 & CIR, $\sqrt{F}$ \\
$A$    & $[0, \infty)$   & Autonomic amplitude (HRV / autonomic resilience)     & CIR, $\sqrt{A + \varepsilon}$ \\
$K_{FB}$ & $[0, \infty)$ & Aerobic fatigue gain (Busso variable-dose)           & CIR, $\sqrt{K_{FB}}$ \\
$K_{FS}$ & $[0, \infty)$ & Strength fatigue gain (Busso variable-dose)          & CIR, $\sqrt{K_{FS}}$ \\
\bottomrule
\end{tabular}
\end{center}

Two operating-point reference constants are baked into the
G1-reparametrised drift: $A_{\rm TYP} = 0.10$ (typical autonomic
amplitude) and $F_{\rm TYP} = 0.20$ (typical fatigue). They calibrate
the model around a physiologically representative state.

\subsection{Drift equations} \label{sec:drift}

The deterministic part of the SDE is implemented in
\path{models/fsa_high_res/_dynamics.py:drift_jax} as the canonical
single source of truth. The six equations are:

\begin{align}
\dot B    &= \kappa_B \, a_B(A) \, \Phi_B - \frac{B}{\tau_B},
   \quad &
   a_B(A) &= \frac{1 + \epsilon_{AB} A}{1 + \epsilon_{AB} A_{\rm TYP}}, \\[3pt]
\dot S    &= \kappa_S \, a_S(A) \, \Phi_S - \frac{S}{\tau_S},
   \quad &
   a_S(A) &= \frac{1 + \epsilon_{AS} A}{1 + \epsilon_{AS} A_{\rm TYP}}, \\[3pt]
\dot F    &= K_{FB}\, \Phi_B + K_{FS}\, \Phi_S - \frac{a_F(A)}{\tau_F}\, F,
   \quad &
   a_F(A) &= \frac{1 + \lambda_A A}{1 + \lambda_A A_{\rm TYP}}, \\[3pt]
\dot A    &= \mu(B, S, F)\,A - \eta\, A^3, \label{eq:dotA}\\[3pt]
\dot K_{FB} &= \frac{K_{FB}^0 - K_{FB}}{\tau_K} + \mu_K\, \Phi_B,
   \quad &
   \dot K_{FS} &= \frac{K_{FS}^0 - K_{FS}}{\tau_K} + \mu_K\, \Phi_S.
\end{align}

The exogenous control input is the bimodal stimulus rate
$\Phi(t) = (\Phi_B(t), \Phi_S(t))^T$ --- aerobic and strength training
loads respectively, both dimensionless.

\paragraph{Stuart--Landau bifurcation parameter (the v5 form).}
The autonomic equation \eqref{eq:dotA} contains the multiplicative
parameter $\mu(B, S, F)$, which is the locus of the v5 deconditioning
extension. Its full v5 form is:
\begin{equation}
\boxed{\;
\mu(B, S, F) \;=\; \mu_0 + \mu_B B + \mu_S S - \mu_F F - \mu_{FF}(F - F_{\rm TYP})^2
   \;-\; \mu_{B-} h_n(B; B_{\rm dec})
   \;-\; \mu_{S-} h_n(S; S_{\rm dec}),
\;}
\label{eq:mu-v5}
\end{equation}
where the v5 Hill-deconditioning weight is
\begin{equation}
h_n(x;\, x_{\rm dec}) \;=\; \frac{x_{\rm dec}^n}{x^n + x_{\rm dec}^n} \;\in\; [0, 1].
\label{eq:hill}
\end{equation}
Setting $\mu_{B-} = \mu_{S-} = 0$ recovers the FSA-v4 form exactly,
which is why $\mu_{\rm dec_*} = 0$ are the v4-recovering defaults in
\texttt{TRUTH\_PARAMS}; \texttt{TRUTH\_PARAMS\_V5} sets them to non-zero
values that activate the closed-island topology of Section
\ref{sec:basins}.

\subsection{What each parameter controls --- physiological intuition}

\begin{center}\small
\begin{tabular}{lp{12cm}}
\toprule
Symbol & Meaning \\
\midrule
$\tau_B, \tau_S$ & Exponential decay time-scales of aerobic and strength capacities under zero stimulus. $\tau_B = 42$d, $\tau_S = 60$d --- so strength is slower to detrain than aerobic fitness. \\
$\kappa_B, \kappa_S$ & Per-unit-stimulus accrual rates for $B$ and $S$, respectively. \\
$\epsilon_{AB}, \epsilon_{AS}$ & Autonomic boost: when $A$ is high, training builds capacity faster ($a_*(A) > 1$). \\
$\tau_F$ & Fatigue clearance time-scale ($\sim 7$d after autonomic correction). \\
$\lambda_A$ & Autonomic acceleration of fatigue clearance: high $A$ clears fatigue faster. \\
$K_{FB}^0, K_{FS}^0$ & Baseline fatigue gains (untrained athlete). $K_{Fi}$ rises above baseline under sustained training. \\
$\tau_K$ & Recovery time-scale of the fatigue sensitivity ($\sim 21$d, three weeks). \\
$\mu_K$ & ``Damage'' rate: how fast training pushes $K$ above baseline. \\
$\mu_0$ & Baseline autonomic drive at zero load. \\
$\mu_B, \mu_S$ & Coupling of fitness and strength to autonomic drive. \\
$\mu_F$ & Fatigue's linear suppression of autonomic drive. \\
$\mu_{FF}$ & Quadratic penalty around $F_{\rm TYP}$ --- captures the cliff where overtraining suppresses HRV. \\
$\eta$ & Cubic damping of $A$ --- prevents runaway autonomic amplitude. \\
\midrule
$B_{\rm dec}, S_{\rm dec}$ & Chronic-fitness thresholds below which deconditioning kicks in. \\
$\mu_{B-}, \mu_{S-}$ & Maximum autonomic-drive penalty when $B$ (resp.\ $S$) is far below threshold. \\
$n_{\rm dec}$ & Hill exponent --- steepness of the deconditioning transition (default 4). \\
\bottomrule
\end{tabular}
\end{center}

\subsection{Diffusion: state-dependent diagonal} \label{sec:diffusion}

The stochastic part of the SDE is diagonal with state-dependent scales,
implemented in \path{_dynamics.py:diffusion_state_dep}. Each component
has a physically-motivated form:
\begin{equation}
\sigma(y) = \mathrm{diag}\!\Big(
  \sigma_B \sqrt{B(1-B)},\;
  \sigma_S \sqrt{S(1-S)},\;
  \sigma_F \sqrt{F},\;
  \sigma_A \sqrt{A + \varepsilon_A},\;
  \sigma_K \sqrt{K_{FB}},\;
  \sigma_K \sqrt{K_{FS}}
\Big).
\label{eq:diffusion}
\end{equation}
The Jacobi-form $\sqrt{x(1-x)}$ scales for $B, S$ keep them strictly in
$[0, 1]$; the CIR-form $\sqrt{x}$ scales for $F, A, K_{*}$ keep them
strictly non-negative. The five $\sigma_*$ scalars and the floor
$\varepsilon_A$ are \emph{frozen} in production --- diffusion-side
identifiability is weak and standard practice is not to estimate
diffusion scales jointly with drift parameters.

\subsection{Cascade structure} \label{sec:cascade}

The drift has a feed-forward cascade structure under \emph{constant}
$\Phi$: $K \to (B, S) \to F \to A$, with the only back-coupling being
the autonomic state $A$ entering the $a_*(A)$ multipliers. This
structure is what makes the v5 stability analysis tractable: at the
slow-manifold equilibrium each forward block has a closed-form
solution given $A$, and the entire 6D fixed-point reduces to a single
scalar self-consistency equation in $A$ (see Section
\ref{sec:basins}).


% =============================================================================
\section{Observation model} \label{sec:obs}
% =============================================================================

The observation likelihood is implemented in
\path{models/fsa_high_res/estimation.py} (functions
\texttt{\_gaussian\_obs\_ll} and \texttt{obs\_log\_weight\_fn}) and the
synthetic-data samplers in \path{simulation.py}
(\texttt{gen\_obs\_sleep, gen\_obs\_hr}, etc.).

\subsection{Five observation channels}

At each observation time $t_k$, with the latent state evaluated at that
bin and a circadian regressor $C_k$, the model emits up to five
observations:

\begin{align}
\mu^{\rm HR}_k  &= {\rm HR}_{\rm base} \;-\; \kappa_B^{\rm HR} B_k \;+\; \alpha_A^{\rm HR} A_k \;+\; \beta_C^{\rm HR} C_k,
   & y^{\rm HR}_k  &\sim \Ncal(\mu^{\rm HR}_k, \sigma_{\rm HR}^2), \label{eq:obs-HR} \\[2pt]
\mu^{\rm S}_k   &= S_{\rm base} \;+\; k_F F_k \;-\; k_{A,S} A_k \;+\; \beta_C^{\rm S} C_k,
   & y^{\rm S}_k   &\sim \Ncal(\mu^{\rm S}_k,  \sigma_S^2), \label{eq:obs-S} \\[2pt]
\mu^{\rm st}_k  &= \mu_{\rm step,0} \;+\; \beta_B^{\rm st} B_k \;-\; \beta_F^{\rm st} F_k \;+\; \beta_A^{\rm st} A_k \;+\; \beta_C^{\rm st} C_k,
   & y^{\rm st}_k  &\sim \Ncal(\mu^{\rm st}_k, \sigma_{\rm st}^2), \label{eq:obs-st} \\[2pt]
\mu^{\rm VL}_k  &= \beta_S^{\rm VL} S_k \;-\; \beta_F^{\rm VL} F_k,
   & y^{\rm VL}_k  &\sim \Ncal(\mu^{\rm VL}_k, \sigma_{\rm VL}^2), \label{eq:obs-VL} \\[2pt]
p^{\rm sleep}_k &= \sigma\!\big(k_C C_k + k_A A_k - \tilde c\big),
   & y^{\rm sleep}_k &\sim \mathrm{Bernoulli}(p^{\rm sleep}_k). \label{eq:obs-sleep}
\end{align}

\paragraph{Linear-in-state.} The four Gaussian channels are linear
functions of the 4-dimensional sub-state $(B, S, F, A)$ plus a circadian
intercept; the $K$ states are not directly observed. The Bernoulli
sleep channel is a logistic in $A$ and $C$. This makes the observation
Jacobian $H$ a constant-coefficient $4 \times 6$ matrix at each time,
enabling the closed-form Kalman fusion in
\path{estimation.py:propagate_fn_v5}.

\paragraph{Reading the channels physiologically.}
\begin{itemize}
  \item \textbf{HR} (heart rate, bpm): increases with autonomic
        amplitude, decreases with chronic aerobic fitness, modulated by
        circadian.
  \item \textbf{Stress} (subjective wellness score): increases with
        fatigue, decreases with autonomic amplitude.
  \item \textbf{Steps} (daily activity log): increases with both
        fitness and autonomic amplitude, decreases with fatigue.
  \item \textbf{Volume Load} (training session volume): observed only
        during training events --- linear in strength capacity, with
        fatigue-driven reduction.
  \item \textbf{Sleep} (binary good/poor): logistic in autonomic
        amplitude and circadian.
\end{itemize}

\subsection{Channel gating: sleep-aware sampling} \label{sec:gating}

Each Gaussian channel is emitted at the bin grid but is \emph{gated} by
the subject's sleep / wake state. Production assumption: a canonical
sleep window of 22:00--06:00 (8 hours), implementable from a single
time-of-day flag. The gating policy is:

\begin{center}\small
\begin{tabular}{lll}
\toprule
Channel & Gating & Rationale \\
\midrule
HR     & sleep-only       & overnight resting HR is the most informative for $B$ \\
Stress & wake-only        & subjective scoring during waking hours \\
Steps  & wake-only        & step counter active only during wake \\
VolumeLoad & training sessions & one entry per training session ($\le 1$/day) \\
Sleep label & every bin    & sleep state itself is the obs \\
\bottomrule
\end{tabular}
\end{center}

The dense overnight HR sampling is not optional: as Section
\ref{sec:integration} discusses, the production 15-minute bin grid
combined with this sleep-gated HR is what makes the model
identifiable. Replacing it with daily-summary HR observations inflates
the FIM condition number by 20 orders of magnitude.

\subsection{Circadian regressor}

The circadian variable $C(t)$ is a single 24-hour cosine,
\begin{equation}
C(t) = \cos(2\pi t + \phi), \qquad \phi := 0,
\end{equation}
implemented in \path{simulation.py:circadian} (numpy) and
\path{simulation.py:circadian_jax} (JAX). The phase $\phi$ is held at
zero in production. The four $\beta_C^*$ coefficients in equations
\eqref{eq:obs-HR}--\eqref{eq:obs-st} together with $k_C$ in
\eqref{eq:obs-sleep} describe the coupling of each channel to this
shared regressor.


% =============================================================================
\section{The three regimes: closed-island basin topology} \label{sec:basins}
% =============================================================================

The headline mathematical content of FSA-v5 is that, for constant
training input $\Phi = (\Phi_B, \Phi_S)$, the long-time behaviour of the
6D system reduces to a single scalar self-consistency equation in $A$,
and the basins of attraction in stimulus space have a closed-island
geometry. This section gives the engineering-level summary; the full
analysis is in \path{main.tex} §7 and §10.

\subsection{Self-consistency equation for $A$}

Under the cascade structure of Section \ref{sec:cascade}, at constant
$\Phi$ each forward block has a closed-form slow-manifold equilibrium
parametrised by $A$:
\begin{align}
K_{Fi}^*(\Phi_i)   &= K_{Fi}^0 + \tau_K \mu_K \Phi_i, \\
B^*(A; \Phi_B)     &= \tau_B \kappa_B \, a_B(A) \, \Phi_B, \\
S^*(A; \Phi_S)     &= \tau_S \kappa_S \, a_S(A) \, \Phi_S, \\
F^*(A; \Phi)       &= \tau_F \, \big(K_{FB}^* \Phi_B + K_{FS}^* \Phi_S\big) / a_F(A).
\end{align}
Substituting these into the v5 bifurcation parameter (eq.\
\ref{eq:mu-v5}) defines the slow-manifold drive $\bar\mu(A; \Phi)$.
The interior equilibria of $A$ then solve
\begin{equation}
\bar\mu(A; \Phi) \;=\; \eta\, A^2,
\label{eq:self-consistency}
\end{equation}
together with the boundary equilibrium $A = 0$. Equation
\eqref{eq:self-consistency} has 0, 1, or 2 positive roots depending on
$\Phi$, giving rise to three qualitatively distinct regimes.

\subsection{The three regimes}

\paragraph{Mono-stable healthy.} $\bar\mu(0; \Phi) > 0$: one positive
root $A^*$. The boundary $A = 0$ is unstable, the interior $A^*$ is
stable. Healthy autonomic amplitude is the unique attractor regardless
of initial condition.

\paragraph{Bistable.} $\bar\mu(0; \Phi) < 0$ but the curve $g(A) =
\bar\mu(A) - \eta A^2$ rises above zero somewhere: two positive roots
$A_{\rm sep} < A^*$. Both $A = 0$ and $A^*$ are stable; $A_{\rm sep}$ is
the unstable separatrix that partitions the basin. Initial $A_0 <
A_{\rm sep}$ collapses to zero, $A_0 > A_{\rm sep}$ rises to $A^*$.

\paragraph{Mono-stable collapsed.} $g$ stays negative on $A > 0$: no
positive roots. $A = 0$ is the global attractor. Autonomic amplitude is
irrecoverable.

\subsection{Closed-island geometry in $(\Phi_B, \Phi_S)$ space}

Figure \ref{fig:closed-island} shows the central qualitative result.
Under the v4 model (left) the healthy region is an open quarter-disc:
sedentary ($\Phi = 0$) is healthy. Under the v5 modification (right)
the boundary is a closed contour strictly interior to both axes: both
$\Phi_B = 0$ (no aerobic) and $\Phi_S = 0$ (no strength) push outside
the healthy region, so deconditioning by either route leads to
collapse.

\begin{figure}[h]
    \centering
    \includegraphics[width=\textwidth]{figures/v5_sedentary_collapse_island.png}
    \caption{Transcritical bifurcation contour $\bar\mu(0;\Phi) = 0$
    under FSA-v4 (left) and FSA-v5 (right). v4 admits a healthy regime
    down to $\Phi = 0$ and along both axes. The v5 Hill-deconditioning
    term ($B_{\rm dec} = S_{\rm dec} = 0.07$, $\mu_{B-} = \mu_{S-} =
    0.10$, $n=4$) closes the contour into a small island whose boundary
    is strictly interior to both coordinate axes. The medical
    ``Goldilocks'' / inverse-U structure emerges as a basin-topology
    change.}
    \label{fig:closed-island}
\end{figure}

Figure \ref{fig:full-bifurcation} adds the saddle-node curve and the
numerical regime classification. Inside the black transcritical
contour: mono-stable healthy. Between black and magenta dashed: bistable
annulus. Outside magenta: mono-stable collapsed.

\begin{figure}[h]
    \centering
    \includegraphics[width=\textwidth]{figures/v5_full_bifurcation.png}
    \caption{FSA-v5 full bifurcation diagram. Left: analytical
    transcritical (black) and saddle-node (magenta dashed) contours.
    Right: numerical regime classification from forward RK4 rollouts
    starting at $A_0 = 0.05$. Bistable annulus surrounds the healthy
    island on all sides; both deconditioning and over-training are
    reachable from a small initial $A$.}
    \label{fig:full-bifurcation}
\end{figure}

\subsection{Time-domain confirmation}

The static analysis is confirmed by deterministic ODE rollouts under
three constant-$\Phi$ scenarios from a common ``trained athlete''
initial state, shown in Figure \ref{fig:trajectories}. \textbf{Both}
sedentary and over-training drive $A \to 0$ but via different
mechanisms and on different time-scales:

\begin{itemize}
  \item \textbf{Sedentary} ($\Phi = 0$): $B$ and $S$ decay
        exponentially at $\tau_B, \tau_S$. Once they cross
        $B_{\rm dec}, S_{\rm dec}$, the Hill terms in \eqref{eq:mu-v5}
        saturate, $\bar\mu$ goes negative, and $A$ collapses around
        day 100--150. Mediated by detraining of chronic capacities.
  \item \textbf{Moderate} ($\Phi = (0.30, 0.30)$, inside the v5 island):
        all states settle to a healthy steady state, $A^* \approx
        0.55$. Goldilocks-band homeostasis.
  \item \textbf{Over-training} ($\Phi = (1.0, 1.0)$): $F$ runs away to
        $\approx 2$ in $\sim 30$ days, $\bar\mu$ goes strongly negative
        via the $-\mu_{FF}(F-F_{\rm TYP})^2$ term, $A$ collapses by day
        $\sim 50$. Mediated by fatigue, not deconditioning.
\end{itemize}

\begin{figure}[h]
    \centering
    \includegraphics[width=\textwidth]{figures/v5_trajectories_three_regimes.png}
    \caption{Deterministic RK4 rollouts under the three regimes.
    Top row: chronic states $B, S, A$. Bottom row: fatigue $F$ and
    sensitivities $K_{FB}, K_{FS}$.}
    \label{fig:trajectories}
\end{figure}

\subsection{Closed-island sanity check}

For verification during a port, the $\bar\mu(0; \Phi)$ values at five
representative stimuli are tabulated below. These should match exactly
when running \texttt{TRUTH\_PARAMS\_V5} through the analytical
$\bar\mu$ function (e.g., \path{control_v5.py:_mu_bar}).

\begin{center}\small
\begin{tabular}{lcl}
\toprule
Stimulus regime & $\bar\mu(0;\Phi)$ & Outcome \\
\midrule
Sedentary,                $\Phi = (0, 0)$    & $-0.180$ & strong collapse \\
Aerobic-only,             $\Phi = (0.30, 0)$ & $-0.059$ & collapse \\
Strength-only,            $\Phi = (0, 0.30)$ & $-0.093$ & collapse \\
Balanced moderate,        $\Phi = (0.30, 0.30)$ & $+0.011$ & \textbf{healthy} ($A^* \approx 0.54$) \\
Over-training,            $\Phi = (1.0, 1.0)$ & $-1.608$ & catastrophic collapse \\
\bottomrule
\end{tabular}
\end{center}


% =============================================================================
\section{The probabilistic cost functional} \label{sec:cost}
% =============================================================================

This section is the v5 main novelty. The basin geometry of Section
\ref{sec:basins} dictates a specific form of cost functional that the
\texttt{smc2fc} controller should optimise.

\subsection{The structurally correct formulation: chance constraints}

Given the bistable annulus surrounding the healthy island, the natural
clinical question is: \emph{what minimum-effort training programme
keeps the athlete in the healthy basin with at least probability $1 -
\alpha$?} This is encoded as a chance constraint on basin escape:
\begin{equation}
\boxed{\;
\Pcal\!\big[\, A_t < A_{\rm sep}(\Phi_t) \,\big] \;\le\; \alpha
   \qquad \forall\, t \in [0, T].
\;}
\label{eq:chance-constraint}
\end{equation}
The cost functional pairs naturally with this constraint:
\begin{equation}
\boxed{\;
\min_\theta \; \lambda_\Phi \int_0^T \|\Phi(t)\|^2 \, \dd t
\quad \text{s.t.}\quad
   \Ecal\!\left[ \int_0^T A_t \, \dd t \right] \;\ge\; A_{\rm target},\;
   \Pcal\!\big[ A_t < A_{\rm sep}(\Phi_t) \big] \;\le\; \alpha.
\;}
\label{eq:v5-min-effort}
\end{equation}
This formulation directly answers the clinical question: minimise total
effort subject to (i) achieving at least target autonomic exposure
$A_{\rm target}$ in expectation and (ii) keeping the basin-escape
probability below the budget $\alpha$ at every time.

\subsection{Why this needs SMC$^2$, not gradient-OT}

The chance constraint \eqref{eq:chance-constraint} is structurally
non-differentiable: it involves an expectation over indicator functions
of state events. A gradient-based open-loop optimiser cannot represent
it; only an SMC$^2$ controller, which carries a particle population
over both parameters and trajectories, can naturally reject or
re-weight $\theta$-particles whose simulated trajectories violate the
constraint at the prescribed rate.

\subsection{Reference implementations: hard / soft variants}

The chance-constraint cost is implemented as \emph{two interchangeable
JIT-friendly variants} in \path{models/fsa_high_res/control_v5.py},
plus a back-compat alias and a legacy NumPy reference. The smc2fc
controller picks one variant based on its outer-loop architecture.

\paragraph{Hard variant — for pure SMC\textsuperscript{2} architectures.}
\texttt{evaluate\_chance\_constrained\_cost\_hard(...)} uses the
indicator $\mathbb{1}[A_t < A_{\rm sep}(\Phi_t)]$ directly. Suitable
when the controller's outer loop weights $\theta$-particles by their
empirical violation rate at each tempering step. The hard indicator
is non-differentiable in $\theta$, so this variant is NOT for HMC inner
kernels. The full implementation is JIT-compiled and \texttt{vmap}-ed
over particles. Per-bin separatrix root-finding is a JAX bisection over
a 64-point grid — no \texttt{scipy.optimize.brentq}, no Python loops.

\paragraph{Soft variant — for HMC architectures.}
\texttt{evaluate\_chance\_constrained\_cost\_soft(...)} replaces the
indicator with a sigmoid surrogate:
\begin{equation}
\mathbb{1}[A_t < A_{\rm sep}(\Phi_t)] \;\longrightarrow\;
\sigma\!\left(\beta \cdot \frac{A_{\rm sep}(\Phi_t) - A_t}{\rm scale}\right),
\label{eq:soft-surrogate}
\end{equation}
which is $C^\infty$ and differentiable through $\theta$ via the trajectory
$A_t$. As $\beta \to \infty$ the surrogate converges to the hard
indicator (the smoke test verifies this within $10^{-3}$ at $\beta =
10^4$). Recommended use: anneal $\beta$ upward over the inference run
— small $\beta$ early on for smooth easy gradients, large $\beta$ later
for structural correctness. Default $\beta = 50$, $\text{scale} = 0.1$
(units of $A$).

\paragraph{Common signature.}

\begin{lstlisting}
def evaluate_chance_constrained_cost_hard(
    theta_particles,        # list of param dicts, OR dict-of-arrays (n_particles,)
    weights,                # (n_particles,), normalised to sum to 1
    Phi_schedule,           # (n_steps, 2)
    *,
    dt: float = 1.0 / 96,   # 15-min bins (FSA_STEP_MINUTES = 15)
    alpha: float = 0.05,    # chance-constraint budget
    A_target: float = 2.0,  # min E[ integral A dt ]
    truth_params_template=None,    # default TRUTH_PARAMS_V5
    initial_state=None,            # (6,) - default = trained athlete
) -> dict:
    # Returns:
    #   mean_effort                  scalar - integral ||Phi||^2 dt
    #   mean_A_integral              scalar - weighted-mean integral A_t dt
    #   violation_rate_per_particle  (n_particles,)
    #   weighted_violation_rate      scalar
    #   satisfies_chance_constraint  bool   - weighted rate <= alpha
    #   satisfies_target             bool   - mean_A_integral >= A_target
    #   A_sep_per_bin                (n_steps,) - analytical separatrix curve
\end{lstlisting}

The soft variant has two extra kwargs, \texttt{beta} (default 50) and
\texttt{scale} (default 0.1); the rest of the signature is identical.

\paragraph{Back-compat alias.}
The original symbol
\texttt{evaluate\_chance\_constrained\_cost} is preserved as an alias
for the hard variant, so existing imports from this module continue to
work unchanged. The slow legacy NumPy implementation lives at
\texttt{\_evaluate\_chance\_constrained\_cost\_legacy} and is kept only
as a debug / readability reference; production callers should use the
hard or soft variants.

\paragraph{Architecture choice for the smc2fc port.}
The decision tree for the smc2fc controller author:
\begin{itemize}
  \item Outer loop is \emph{pure SMC\textsuperscript{2} with importance
        weighting} — particles are re-weighted at each tempering step
        by their per-particle violation rate, no inner gradient kernel.
        Use \textbf{hard} variant.
  \item Outer loop has \emph{HMC} (or any gradient-based optimiser)
        inside the SMC tempering — needs $\nabla_\theta J$. Use
        \textbf{soft} variant; anneal $\beta$ from $\sim 10$ at the
        first temperature to $\sim 10^3$ near convergence.
  \item Need a back-compat baseline for unit tests against the v3 cost
        functional — use \texttt{build\_control\_spec\_v5}
        (gradient-OT) only as a fallback.
\end{itemize}

The smc2fc controller is responsible for closing the loop:
re-weighting, rejecting, or scoring candidate $\Phi$ schedules using
the dict returned from either variant. This file does not own the
controller; it provides only the cost evaluators.

\subsection{Engineering note: deterministic CRN}

Both variants forward-simulate each particle deterministically (no
per-particle noise) for speed and reproducibility. The chance
constraint then captures \emph{parameter uncertainty}, not state-noise
uncertainty. To upgrade to full stochastic simulation, replace the
internal helper \texttt{\_make\_forward\_rollout\_fn} with a CRN-noise
variant — every other piece of the cost function logic is unchanged.

\subsection{Bonus formulation: Lagrangian relaxation}

For consumers that need a smooth scalar cost (e.g.\ for warmstart of
the SMC$^2$ controller), the chance constraint can be relaxed to a
soft penalty:
\begin{equation}
\widetilde J(\theta) \;=\; \lambda_\Phi \int \|\Phi\|^2 \, \dd t
   \;-\; \Ecal\!\left[\int A \, \dd t\right]
   \;+\; \lambda_{\rm chance} \, \max\!\big(0,\; \widehat p_{\rm violate} - \alpha\big)^2,
\end{equation}
where $\widehat p_{\rm violate}$ is the empirical violation rate from
the particle cloud and $\lambda_{\rm chance}$ is a large penalty
weight. This loses the chance-constraint guarantee but yields a smooth
objective amenable to gradient methods. Equivalent to using the soft
variant with $\beta \to 0$ (very loose surrogate) plus an explicit
violation-rate penalty.


% =============================================================================
\section{Numerical integration: time step and stiffness} \label{sec:integration}
% =============================================================================

\subsection{Production grid: 15-minute bins}

The model is integrated on a 15-minute bin grid:
\begin{equation}
\Delta t \;=\; \frac{1}{96}\;\text{day} \;\approx\; 15\;\text{minutes}, \qquad
\text{bins per day} \;=\; 96.
\end{equation}
This grid is set by the environment variable \texttt{FSA\_STEP\_MINUTES}
(default 15) in \path{models/fsa_high_res/simulation.py}. The
computation is:
\begin{lstlisting}
import os as _os
_STEP_MIN = int(_os.environ.get('FSA_STEP_MINUTES', '15'))
BINS_PER_DAY = (60 * 24) // _STEP_MIN          # = 96 at default
DT_BIN_DAYS  = 1.0 / BINS_PER_DAY              # = 1/96
\end{lstlisting}

Each integration step is exactly one bin --- there is no inner
sub-stepping in the production path
(\path{_plant.py:_plant_em_step}). The choice is justified by stiffness
analysis below.

\subsection{Why one EM step per bin is safe}

The slowest characteristic time is $\tau_S = 60$d (strength capacity
relaxation). The fastest is $\tau_F = 6.36$d (fatigue clearance, after
the autonomic correction). The $A$-equation has a cubic damping term
$-\eta A^3$ which can be locally fast for large $A$, but $A \in [0,
\sim 1]$ in practice and $\eta = 0.20$, so the autonomic time constant
is no faster than $\sim 5$d.

The shortest meaningful time-scale is therefore $\sim 6$d. The
integration step is $\Delta t = 1/96 \approx 0.01$d --- about
$600\times$ smaller than the fastest mode. Forward Euler--Maruyama is
unconditionally stable on this margin; no implicit or sub-stepped
integrator is needed. (For tools that wish to use coarser dt --- e.g.\
the basin sweep tool --- a sub-stepped helper is provided as
\path{_dynamics.py:imex_step_substepped}.)

\subsection{Why daily sampling is non-viable for inference}

Sampling resolution is non-negotiable. Replacing the production
15-minute observation grid with daily summaries inflates the FIM
condition number from $\sim 10^{18}$ to $\sim 10^{38}$ and creates an
artificial $(\epsilon_{AS}, \mu_S)$ confound that the production
sampling does not have. Two specific reasons:

\begin{enumerate}
  \item Dense overnight resting-HR sampling pins down aerobic fitness
        $B$ from passive autonomic signals --- a much more informative
        window than sparse daytime spot-checks.
  \item Time-resolved sleep-gated HR separates the
        $(\epsilon_{AS}, \mu_S)$ contributions to autonomic
        amplitude that are otherwise confounded.
\end{enumerate}

\textbf{Recommendation for the smc2fc port: keep \texttt{FSA\_STEP\_MINUTES = 15}}
and accept the higher computational cost (8064 SDE steps over 84 days
versus $\sim 340$ at $\Delta t = 0.25$d).

\subsection{State clipping after each step}

The post-Euler--Maruyama state is clipped back into the physical
domain at each bin:

\begin{itemize}
  \item $B \in [\varepsilon_B,\; 1 - \varepsilon_B]$ (Jacobi reflection or clip)
  \item $S \in [\varepsilon_S,\; 1 - \varepsilon_S]$ (same)
  \item $F, A, K_{FB}, K_{FS} \in [0, \infty)$ (\texttt{max(\(\cdot\), 0)})
\end{itemize}

with $\varepsilon_B = \varepsilon_S = \varepsilon_A = 10^{-4}$. The
reflection on $B, S$ matches the natural behaviour of the Jacobi
diffusion process at boundaries; the half-line clipping for $F, A, K$
is a safety net since the CIR diffusion is degenerate at zero.

\subsection{Numerical convention summary}

\begin{center}\small
\begin{tabular}{ll}
\toprule
Quantity & Value (production) \\
\midrule
Time unit                 & days \\
Bin width $\Delta t$      & $1/96$ d $\approx 15$ min \\
Bins per day              & $96$ \\
Sub-stepping in plant     & none (single EM step per bin) \\
Sub-stepping in tools     & 4 sub-steps optional via \texttt{imex\_step\_substepped} \\
Reflection: $B, S$        & $\varepsilon = 10^{-4}$ on both ends \\
Floor: $F, A, K_*$        & $0$, with $A$ getting a $+\varepsilon_A$ floor in diffusion only \\
Random number generator   & JAX (CRN-friendly via \texttt{jax.random.PRNGKey}) \\
Float precision           & enable via \texttt{JAX\_ENABLE\_X64=True} (Section \ref{sec:bugs}) \\
\bottomrule
\end{tabular}
\end{center}


% =============================================================================
\section{Parameter set: 37 estimated, 14 frozen} \label{sec:pinning}
% =============================================================================

This section explains the pinning strategy used in
\path{HIGH_RES_FSA_V5_ESTIMATION}. Read the appendix for the full
symbol-by-symbol table; this section gives the rationale.

\subsection{The split}

\begin{center}
\begin{tabular}{ll}
\toprule
Estimated (37) & in \texttt{PARAM\_PRIOR\_CONFIG} \\
Frozen (14)    & in \texttt{frozen\_params} \\
\bottomrule
\end{tabular}
\end{center}

The 14 frozen parameters fall into two groups:

\begin{enumerate}
  \item \textbf{Diffusion-side (6 always frozen):}
        $\sigma_B, \sigma_S, \sigma_F, \sigma_A, \sigma_K, \phi$.
        Diffusion identifiability is weak under SMC$^2$ in any setting
        and is not the right thing to learn jointly with drift.
  \item \textbf{Dynamics-side (8 frozen by FIM analysis + first-port-test policy):}
        $K_{FB}^0,\, K_{FS}^0,\, \tau_K,\, n_{\rm dec},\, B_{\rm dec},\,
        S_{\rm dec},\, \mu_{B-},\, \mu_{S-}$.
\end{enumerate}

\subsection{Why each dynamics-side freeze}

\paragraph{$K_{FB}^0$ and $K_{FS}^0$ (structurally non-identifiable).}
The $K$ subsystem rapidly relaxes onto its slow manifold $K_{Fi}^* =
K_{Fi}^0 + \tau_K \mu_K \Phi_i$. The data observes $K_{Fi}$ only
indirectly through the $F$-equation, and only ever sees the
\emph{combination} $K_{Fi}^0 + \tau_K \mu_K \Phi_i$ --- never $K_{Fi}^0$
in isolation. No amount of additional data fixes this: the FIM
eigenvector for these is structurally zero. Pin them to physiologically
sensible values ($K_{FB}^0 = 0.030$, $K_{FS}^0 = 0.050$) and forget
about estimating them.

\paragraph{$\tau_K$ (jointly weakly identified with $\mu_K$).}
At steady state, only the product $\tau_K \mu_K$ is observable from $K$
relaxation. To separate them you need calibration data with sustained
loading periods of order $\tau_K \approx 21$ days, which is rarely
available. Pin $\tau_K$ to the Busso-standard $21$d and estimate only
$\mu_K$.

\paragraph{$n_{\rm dec}$ (structural shape).}
The Hill exponent is a structural shape parameter, not a data-fitted
one. The transition steepness it controls is a modelling choice. Pin
at $n_{\rm dec} = 4$.

\paragraph{$B_{\rm dec}, S_{\rm dec}, \mu_{B-}, \mu_{S-}$ (need long detraining episodes).}
The four v5 deconditioning parameters are only informable from data
that includes long ($\ge 6$ weeks) detraining episodes, because the
Hill term in equation \eqref{eq:mu-v5} is silent above the threshold.
For a first port-test on simple steady-state synthetic data, pin them
all at the §10 Figure 7 calibration values:
$(B_{\rm dec}, S_{\rm dec}, \mu_{B-}, \mu_{S-}) = (0.07, 0.07, 0.10, 0.10)$.

\subsection{Relaxation policy: how to widen the parameter set later}

When richer calibration data becomes available, freezes can be relaxed
one at a time by moving a key from \texttt{frozen\_params} into
\texttt{PARAM\_PRIOR\_CONFIG}. Suggested priorities, in order of
unblockability:

\begin{enumerate}
  \item $\mu_{B-}, \mu_{S-}$ if a long detraining episode exists in the data.
  \item $B_{\rm dec}, S_{\rm dec}$ if the detraining episode crosses the threshold smoothly.
  \item $\tau_K$ if a $\ge 21$d sustained-load episode is available.
  \item $K_{FB}^0, K_{FS}^0$ --- never. Structurally non-identifiable.
\end{enumerate}

\subsection{Quick reference: what's in the parameter vector}

The \texttt{HIGH\_RES\_FSA\_V5\_ESTIMATION} parameter vector is 37 long.
Component-wise view:

\begin{center}\small
\begin{tabular}{lll}
\toprule
Group & Count & Members (briefly) \\
\midrule
Aerobic-capacity dynamics  & 3  & $\tau_B, \kappa_B, \epsilon_{AB}$ \\
Strength-capacity dynamics & 3  & $\tau_S, \kappa_S, \epsilon_{AS}$ \\
Fatigue dynamics           & 2  & $\tau_F, \lambda_A$ \\
Variable-dose K dynamics   & 1  & $\mu_K$ ($\tau_K$ frozen, $K^0_{Fi}$ frozen) \\
Stuart--Landau drive       & 6  & $\mu_0, \mu_B, \mu_S, \mu_F, \mu_{FF}, \eta$ \\
HR observation             & 5  & ${\rm HR}_{\rm base}, \kappa_B^{\rm HR}, \alpha_A^{\rm HR}, \beta_C^{\rm HR}, \sigma_{\rm HR}$ \\
Sleep Bernoulli            & 3  & $k_C, k_A, \tilde c$ \\
Stress observation         & 5  & $S_{\rm base}, k_F, k_{A,S}, \beta_C^{\rm S}, \sigma_S$ \\
Steps observation          & 6  & $\mu_{\rm step,0}, \beta_B^{\rm st}, \beta_F^{\rm st}, \beta_A^{\rm st}, \beta_C^{\rm st}, \sigma_{\rm st}$ \\
Volume Load observation    & 3  & $\beta_S^{\rm VL}, \beta_F^{\rm VL}, \sigma_{\rm VL}$ \\
\midrule
\textbf{Total estimated}   & \textbf{37} & \\
\bottomrule
\end{tabular}
\end{center}


% =============================================================================
\section{Test scenarios for synthetic data} \label{sec:tests}
% =============================================================================

This section gives concrete test scenarios that exercise the v5
implementation end-to-end. Each one is a constant-$\Phi$ rollout from a
defined initial state, with a predicted qualitative outcome that the
engineer can verify against. The test infrastructure is in
\path{tests/test_fsa_v5_smoke.py}.

\subsection{Test 1: Sedentary collapse}

\textbf{Setup:} Trained-athlete initial state $(B, S, F, A, K_{FB},
K_{FS}) = (0.50, 0.45, 0.20, 0.45, 0.06, 0.07)$. Constant
$\Phi = (0, 0)$. Run for 200 days.

\textbf{Expected:} $B$ and $S$ decay exponentially at $\tau_B = 42$d
and $\tau_S = 60$d. Around day 60--100, both fall below the
deconditioning threshold $B_{\rm dec} = S_{\rm dec} = 0.07$ and the
Hill terms in $\mu$ saturate. With $\bar\mu$ now strongly negative,
$A$ collapses to near zero by day $\sim 150$. Fatigue $F$ also decays
(no input). $K$ states relax back to baseline.

\textbf{Pass criterion:} $A(t = 200\text{d}) < 0.01$, but $A(t =
30\text{d}) > 0.3$ (not collapsed yet) --- demonstrates the
slow-detraining time-course.

\textbf{Code:}
\begin{lstlisting}
from models.fsa_high_res import StepwisePlant, DEFAULT_PARAMS_V5
import numpy as np

plant = StepwisePlant(truth_params=dict(DEFAULT_PARAMS_V5))
plant.state = np.array([0.50, 0.45, 0.20, 0.45, 0.06, 0.07])
out = plant.advance(stride_bins=200 * 96,
                     Phi_daily=np.zeros((200, 2)))
A_traj = out['trajectory'][:, 3]
assert A_traj[-1] < 0.01,    f"A did not collapse: {A_traj[-1]}"
assert A_traj[30 * 96] > 0.3, f"A collapsed too early"
\end{lstlisting}

\subsection{Test 2: Moderate (healthy) attractor}

\textbf{Setup:} Same trained-athlete initial state. Constant
$\Phi = (0.30, 0.30)$ (inside the v5 healthy island). Run for 100 days.

\textbf{Expected:} $B$ and $S$ relax to slow-manifold values near
$0.18$ each. $\bar\mu$ stays positive ($\sim +0.011$). $A$ relaxes to
its healthy attractor $A^* \approx 0.55$ over the first $\sim 50$ days.
$F$ stabilises near $F_{\rm TYP} = 0.20$. $K$ states rise modestly.

\textbf{Pass criterion:} $A(100\text{d}) > 0.4$, $|F(100\text{d}) - 0.20| < 0.10$, no
NaN.

\subsection{Test 3: Over-training collapse}

\textbf{Setup:} Same trained-athlete initial state. Constant $\Phi =
(1.0, 1.0)$ (well outside the v5 island). Run for 60 days.

\textbf{Expected:} $F$ runs away to $\approx 2$ within 30 days. The
$-\mu_{FF}(F-F_{\rm TYP})^2$ term in $\mu$ dominates, $\bar\mu$ goes
strongly negative ($\sim -1.6$), and $A$ collapses to zero by day
$\sim 50$. $B, S$ saturate at their upper bounds.

\textbf{Pass criterion:} $A(60\text{d}) < 0.01$, but $F(60\text{d}) >
1.5$ (the collapse mechanism is fatigue-driven, not deconditioning).

\subsection{Test 4: Bistability sensitivity (separatrix crossing)}

\textbf{Setup:} Constant $\Phi$ chosen inside the bistable annulus
(e.g.\ $\Phi = (0.50, 0.30)$ from §7 of \path{main.tex}). Run from two
different initial states differing only in $A_0$:
\begin{itemize}
  \item Trial A: $A_0 = 0.10$ (above separatrix)
  \item Trial B: $A_0 = 0.02$ (below separatrix)
\end{itemize}

\textbf{Expected:} Trial A converges to the healthy interior
attractor; Trial B collapses to $A = 0$. Same training, different
outcomes --- this is the bistability test.

\textbf{Pass criterion:} $|A_{\rm trial A}(150\text{d}) - 0.5| < 0.2$ and
$A_{\rm trial B}(150\text{d}) < 0.01$.

\subsection{Test 5: Closed-loop with the chance-constrained cost}

\textbf{Setup:} Build a 10-particle cloud of parameters all at the v5
truth (no inference yet). Build a 14-day moderate $\Phi = (0.30, 0.30)$
schedule. Call \texttt{evaluate\_chance\_constrained\_cost} with
$\alpha = 0.05$, $A_{\rm target} = 2.0$.

\textbf{Expected:} \texttt{satisfies\_chance\_constraint = True},
\texttt{satisfies\_target = True}, \texttt{weighted\_violation\_rate}
$\approx 0$, \texttt{mean\_A\_integral} $\approx 11$ ($\approx
A^* \cdot 14\text{d} \cdot 96$ bins$\cdot$ dt).

This test exercises the entire v5 cost-evaluation path end-to-end.

\subsection{Recommended order to run during the port}

\begin{enumerate}
  \item Test 2 (moderate, healthy) --- confirms the implementation can produce a stable trajectory.
  \item Test 1 (sedentary collapse) --- confirms the Hill term fires correctly under low $\Phi$.
  \item Test 3 (over-training) --- confirms the v4-inherited fatigue mechanism is intact.
  \item Test 4 (bistability) --- confirms the closed-island geometry actually exists at runtime.
  \item Test 5 (cost evaluator) --- confirms the cost evaluator wires up correctly.
\end{enumerate}

If any of these fails, the port has an issue; consult Section
\ref{sec:bugs}.


% =============================================================================
\section{Common bugs and gotchas} \label{sec:bugs}
% =============================================================================

This section documents specific issues encountered during the v5
implementation pass. Read it before you port; each item below is a
real bug that took non-trivial effort to surface.

\subsection{The \texttt{sigma\_S} name collision}

\textbf{Symptom:} The strength state $S$ in the latent SDE develops
500$\times$ larger noise than expected; trajectory shows wild swings of
order $\pm 0.5$ over a single day.

\textbf{Cause:} The \texttt{DEFAULT\_PARAMS} dictionary in
\path{simulation.py} contains \emph{two} keys named \texttt{sigma\_S} ---
one for the latent-S Jacobi diffusion scale ($\sim 0.008$) and another
for the stress-channel observation noise ($\sim 4.0$). Python's dict
keeps only the last assignment, which is the observation one. Reading
\texttt{truth\_params['sigma\_S']} therefore returns $4.0$ (the obs
noise) where the SDE integrator expects the diffusion scale.

\textbf{Fix:} The plant in \path{_plant.py} hard-codes the diffusion
scales as constants matching \texttt{SIGMA\_*\_FROZEN} in
\path{estimation.py}, bypassing the params dict entirely. The
\path{estimation.py} \texttt{propagate\_fn} also hard-codes them. If
you write new code that needs the diffusion scales, follow this
convention --- do not read them from the params dict.

\subsection{\texttt{JAX\_ENABLE\_X64} required}

\textbf{Symptom:} ``Explicitly requested dtype float64 ... is not
available, and will be truncated to dtype float32'' warnings spam the
console; numerical results subtly disagree with reference values.

\textbf{Cause:} JAX defaults to float32 unless either the environment
variable \texttt{JAX\_ENABLE\_X64=True} is set or
\texttt{jax.config.update("jax\_enable\_x64", True)} is called before
JAX is imported.

\textbf{Fix:} Set the environment variable, or do
\begin{lstlisting}
import os
os.environ['JAX_ENABLE_X64'] = 'True'
from jax import config
config.update('jax_enable_x64', True)
\end{lstlisting}
at the top of any tool/test that runs the model. The smoke test
(\path{tests/test_fsa_v5_smoke.py}) does this.

\subsection{\texttt{EstimationModel} requires \texttt{init\_state\_prior\_config}}

\textbf{Symptom:} \texttt{TypeError: EstimationModel.\_\_init\_\_() missing
1 required positional argument: 'init\_state\_prior\_config'}.

\textbf{Cause:} The \texttt{EstimationModel} dataclass in
\texttt{smc2fc/estimation\_model.py} has
\texttt{init\_state\_prior\_config} as a required (not optional)
constructor argument.

\textbf{Fix:} Pass an empty \texttt{OrderedDict()} for it if no
init-state prior is desired. The v5 \texttt{HIGH\_RES\_FSA\_V5\_ESTIMATION}
in \path{estimation.py} does exactly this.

\subsection{Drift duplication --- single source of truth}

\textbf{Symptom:} Subtle drift differences between the JAX path and
the numpy path that only manifest at certain parameter regimes.

\textbf{Cause:} The drift function appears in three places:
\path{_dynamics.py:drift_jax} (canonical), \path{simulation.py:drift}
(numpy mirror), and \path{simulation.py:drift_jax} (JAX wrapper for
SDEModel API). All three must be kept in sync. The v5 implementation
delegates from \path{estimation.py} to the canonical
\path{_dynamics.drift_jax} via factory pattern, but the
\path{simulation.py} copies are still inlined (because they have a
different calling convention).

\textbf{Fix:} If you modify the drift, modify all three places. There
is a comment block at the top of \path{simulation.py:drift} that says
``KEEP IN SYNC with \texttt{\_dynamics.drift\_jax}''. Read that. The
chance-constraint cost in \path{control_v5.py} also re-derives
$\bar\mu$ in a fourth place --- keep that one in sync too.

\subsection{Plant truth\_params must include observation coefficients}

\textbf{Symptom:} \texttt{KeyError: 'k\_C'} (or any obs coefficient)
when running \texttt{StepwisePlant.advance()}.

\textbf{Cause:} The plant calls the obs samplers
(\texttt{gen\_obs\_sleep, gen\_obs\_hr, ...}), which read observation
coefficients from \texttt{self.truth\_params}. If you initialise the
plant with \texttt{truth\_params=TRUTH\_PARAMS\_V5} (from
\path{_dynamics.py}), it has only the dynamics + diffusion keys ---
not the obs coefficients.

\textbf{Fix:} Initialise the plant with
\texttt{truth\_params=DEFAULT\_PARAMS\_V5} (from \path{simulation.py}),
which has both. Or merge: \texttt{dict(\{**DEFAULT\_PARAMS\_V5,
**TRUTH\_PARAMS\_V5\})}.

\subsection{$K_{FB}^0$ and $K_{FS}^0$ are structurally non-identifiable}

\textbf{Symptom:} The posterior on $K_{FB}^0$ and $K_{FS}^0$ looks
identical to the prior; condition number of the FIM blows up.

\textbf{Cause:} As Section \ref{sec:pinning} explains, the data only
ever sees the combination $K_{Fi}^0 + \tau_K \mu_K \Phi_i$. The
constants $K_{Fi}^0$ are formally non-identifiable.

\textbf{Fix:} Pin them. The default v5 \texttt{frozen\_params} dict
already does this (see \path{estimation.py}). Do not be tempted to add
them to \texttt{PARAM\_PRIOR\_CONFIG} no matter how rich the data is.

\subsection{Don't collapse observations to daily summaries during inference}

\textbf{Symptom:} Suspiciously bad posterior identifiability;
$(\epsilon_{AS}, \mu_S)$ pair always shows up in the FIM smallest
eigenvectors.

\textbf{Cause:} Switching from the production 15-minute bin grid to
daily summary observations destroys the dense overnight HR signal that
pins $B$ and the time-resolved sleep-gated HR that separates
$\epsilon_{AS}$ from $\mu_S$ (Section \ref{sec:integration}).

\textbf{Fix:} Keep \texttt{FSA\_STEP\_MINUTES = 15}. Pay the
$\sim 24\times$ extra computational cost. The FIM analysis (§11 of
\path{main.tex}) shows this changes the condition number by 20 orders
of magnitude.

\subsection{Particle parameter dicts must include v5 Hill keys}

\textbf{Symptom:} \texttt{KeyError: 'B\_dec'} (or any v5 Hill key) when
forward-simulating particles.

\textbf{Cause:} The canonical \texttt{drift\_jax} in
\path{_dynamics.py} expects all v5 Hill keys present (\texttt{B\_dec,
S\_dec, mu\_dec\_B, mu\_dec\_S, n\_dec}). If the SMC$^2$ engine carries
$\theta$-particles whose dicts only contain the estimated parameters
(37 of them, Section \ref{sec:pinning}), the four v5 Hill keys are
missing and the drift call fails.

\textbf{Fix:} Always merge with the frozen-params dict before calling
\texttt{drift\_jax}. The cost evaluator
\texttt{evaluate\_chance\_constrained\_cost} in
\path{control_v5.py} already does this; new consumers need to follow
suit. Pattern:
\begin{lstlisting}
p_full = {**FROZEN_DICT, **estimated_dict}
y_next = y + dt * drift_jax(y, p_full, Phi_t)
\end{lstlisting}

\subsection{Observation gating must match production}

\textbf{Symptom:} FIM analysis under-estimates information content;
inference posterior is wider than expected.

\textbf{Cause:} The four Gaussian observation channels are gated by
sleep/wake state (HR sleep-only, Stress and Steps wake-only, Volume
Load training-session-only). If your inference pipeline ingests the
observations as if they were emitted at every bin, the Kalman fusion
weights them all and the FIM looks fatter than it really is. The
opposite mistake (treating them as daily averages) destroys
identifiability as in §\ref{sec:integration}.

\textbf{Fix:} Use the \texttt{align\_obs\_fn} in
\path{estimation.py:align_obs_fn} to pack scattered observations into
bin-resolved arrays with explicit \texttt{*\_present} masks. The
\texttt{propagate\_fn\_v5} respects the masks correctly.

\subsection{Initial state defaults: bounded but ``deconditioned''}

The default initial state in \texttt{DEFAULT\_INIT}
($B = 0.05, S = 0.10, F = 0.30, A = 0.10, K_{FB} = 0.030, K_{FS} =
0.050$) is a ``deconditioned but otherwise healthy'' starting point,
not a trained athlete. Do not be surprised if a forward sim from
\texttt{DEFAULT\_INIT} under moderate $\Phi$ takes weeks to settle ---
the system is climbing from low capacity. For testing closed-loop
behaviour, override the initial state with a trained-athlete vector
(Section \ref{sec:tests}).


% =============================================================================
\section*{Appendix: Complete FSA-v5 Parameter Reference}
\addcontentsline{toc}{section}{Appendix: Complete FSA-v5 Parameter Reference}
\label{sec:appendix-v5-parameters}

This appendix is the canonical reference for every parameter in the
FSA-v5 model as implemented in the ``smc2fc``-port-ready codebase. For
each parameter we list its symbol, role in the equations of motion,
prior (if estimated) or frozen value, and where it first appears in the
main document.

The parameter set as implemented in
\path{models/fsa_high_res/estimation.py} consists of \textbf{37
estimated parameters} plus \textbf{14 frozen} (8 dynamics +
6 diffusion-related). The maximum-pinning rationale is the §11.6
recommendation derived from the Fisher-information analysis in §11.

\subsection*{Estimated parameters (37)}

These are the entries of \texttt{PARAM\_PRIOR\_CONFIG} in
\texttt{estimation.py}. Lognormal priors are parameterised as
$\log\theta \sim \mathcal{N}(\mu_{\log}, \sigma_{\log}^2)$ — the table
shows $(\exp(\mu_{\log}),\, \sigma_{\log})$. Normal priors are
parameterised directly as $(\mu, \sigma)$.

\begin{center}\small
\begin{tabular}{llll}
\toprule
Symbol & Role / location & Prior & Truth \\
\midrule
\multicolumn{4}{l}{\textit{Dynamics — aerobic capacity B}} \\
\midrule
$\tau_B$ & B time-scale / §11.1 & lognormal $(42.0,\,0.10)$ & $42.0$ d \\
$\kappa_B$ & B gain / §11.1 & lognormal $(0.01248,\,0.20)$ & $0.01248$ \\
$\epsilon_{AB}$ & autonomic-coupling on B / §11.1 & lognormal $(0.40,\,0.05)$ & $0.40$ \\
\midrule
\multicolumn{4}{l}{\textit{Dynamics — strength capacity S}} \\
\midrule
$\tau_S$ & S time-scale / §11.1 & lognormal $(60.0,\,0.10)$ & $60.0$ d \\
$\kappa_S$ & S gain / §11.1 & lognormal $(0.00816,\,0.20)$ & $0.00816$ \\
$\epsilon_{AS}$ & autonomic-coupling on S / §11.1 & lognormal $(0.20,\,0.05)$ & $0.20$ \\
\midrule
\multicolumn{4}{l}{\textit{Dynamics — fatigue F}} \\
\midrule
$\tau_F$ & F time-scale / §11.1 & lognormal $(6.36,\,0.15)$ & $6.36$ d \\
$\lambda_A$ & autonomic-coupling on F / §11.1 & lognormal $(1.00,\,0.05)$ & $1.00$ \\
\midrule
\multicolumn{4}{l}{\textit{Dynamics — variable-dose K (Busso 2003)}} \\
\midrule
$\mu_K$ & K damage rate / §11.1 & lognormal $(0.005,\,0.20)$ & $0.005$ \\
\midrule
\multicolumn{4}{l}{\textit{Dynamics — Stuart--Landau bifurcation parameter $\bar\mu$}} \\
\midrule
$\mu_0$ & baseline drive / §11.1 & lognormal $(0.036,\,0.20)$ & $0.036$ \\
$\mu_B$ & B$\to$A coupling / §11.1 & lognormal $(0.30,\,0.20)$ & $0.30$ \\
$\mu_S$ & S$\to$A coupling / §11.1 & lognormal $(0.15,\,0.20)$ & $0.15$ \\
$\mu_F$ & F$\to$A linear suppression / §11.1 & lognormal $(0.26,\,0.20)$ & $0.26$ \\
$\mu_{FF}$ & quadratic F penalty / §11.1 & lognormal $(0.40,\,0.05)$ & $0.40$ \\
$\eta$ & cubic damping in A / §11.1 & lognormal $(0.20,\,0.15)$ & $0.20$ \\
\midrule
\multicolumn{4}{l}{\textit{HR observation channel}} \\
\midrule
$\text{HR}_\text{base}$ & resting HR baseline (bpm) / §11.1 & normal $(62.0,\,2.0)$ & $62.0$ \\
$\kappa_B^{HR}$ & HR$\propto -B$ slope / §11.1 & lognormal $(12.0,\,0.15)$ & $12.0$ \\
$\alpha_A^{HR}$ & HR$\propto +A$ slope / §11.1 & lognormal $(3.0,\,0.20)$ & $3.0$ \\
$\beta_C^{HR}$ & HR circadian coefficient / §11.1 & normal $(-2.5,\,0.5)$ & $-2.5$ \\
$\sigma_{HR}$ & HR observation noise (bpm) / §11.1 & lognormal $(2.0,\,0.20)$ & $2.0$ \\
\midrule
\multicolumn{4}{l}{\textit{Sleep Bernoulli channel}} \\
\midrule
$k_C$ & sleep$\propto C$ slope / §11.1 & lognormal $(3.0,\,0.15)$ & $3.0$ \\
$k_A$ & sleep$\propto A$ slope / §11.1 & lognormal $(2.0,\,0.25)$ & $2.0$ \\
$\tilde c$ & sleep logit intercept / §11.1 & normal $(0.5,\,0.25)$ & $0.5$ \\
\midrule
\multicolumn{4}{l}{\textit{Stress observation channel}} \\
\midrule
$S_\text{base}$ & stress baseline / §11.1 & normal $(30.0,\,3.0)$ & $30.0$ \\
$k_F$ & stress$\propto +F$ slope / §11.1 & lognormal $(20.0,\,0.20)$ & $20.0$ \\
$k_{A,S}$ & stress$\propto -A$ slope / §11.1 & lognormal $(8.0,\,0.25)$ & $8.0$ \\
$\beta_C^{S}$ & stress circadian coefficient / §11.1 & normal $(-4.0,\,0.8)$ & $-4.0$ \\
$\sigma_S$ & stress observation noise / §11.1 & lognormal $(4.0,\,0.20)$ & $4.0$ \\
\midrule
\multicolumn{4}{l}{\textit{Steps observation channel (log-Gaussian)}} \\
\midrule
$\mu_\text{step,0}$ & log-steps intercept / §11.1 & normal $(5.5,\,0.3)$ & $5.5$ \\
$\beta_B^{st}$ & steps$\propto +B$ slope / §11.1 & lognormal $(0.8,\,0.20)$ & $0.8$ \\
$\beta_F^{st}$ & steps$\propto -F$ slope / §11.1 & lognormal $(0.5,\,0.25)$ & $0.5$ \\
$\beta_A^{st}$ & steps$\propto +A$ slope / §11.1 & lognormal $(0.3,\,0.25)$ & $0.3$ \\
$\beta_C^{st}$ & steps circadian coefficient / §11.1 & normal $(-0.8,\,0.2)$ & $-0.8$ \\
$\sigma_\text{st}$ & log-steps noise / §11.1 & lognormal $(0.5,\,0.15)$ & $0.5$ \\
\midrule
\multicolumn{4}{l}{\textit{Volume Load observation channel}} \\
\midrule
$\beta_S^{VL}$ & VL$\propto +S$ slope / §11.1 & lognormal $(100.0,\,0.15)$ & $100.0$ \\
$\beta_F^{VL}$ & VL$\propto -F$ slope / §11.1 & lognormal $(20.0,\,0.20)$ & $20.0$ \\
$\sigma_{VL}$ & VL observation noise / §11.1 & lognormal $(10.0,\,0.20)$ & $10.0$ \\
\bottomrule
\end{tabular}
\end{center}

\subsection*{Frozen parameters (8 dynamics + 6 diffusion-related)}

These do not appear in the estimated parameter vector. They are baked
into the ``EstimationModel.frozen\_params`` dict. Reasons for freezing
each one are tied directly to the §11 FIM analysis.

\begin{center}\small
\begin{tabular}{lllp{6cm}}
\toprule
Symbol & Frozen value & Role / location & Reason for freezing \\
\midrule
\multicolumn{4}{l}{\textit{Dynamics-side freezes}} \\
\midrule
$K_{FB}^0$  & $0.030$ & baseline aerobic K / §11.1   & Structurally non-identifiable. Slow-manifold relaxation collapses $K_{FB}^0$ and $\tau_K \mu_K$ into a single observable. §11.6 result (1). \\
$K_{FS}^0$  & $0.050$ & baseline strength K / §11.1  & Same as $K_{FB}^0$. \\
$\tau_K$    & $21.0$ d & K relaxation timescale / §11.1 & Busso-standard. Jointly weakly identified with $\mu_K$ (§11.6 result 5). \\
$n_\text{dec}$ & $4.0$ & Hill exponent / §10 eq.\ 25 & Structural shape parameter; not learned. \\
$B_\text{dec}$ & $0.07$ & B deconditioning threshold / §10 & Pinned to §10.4 Table 5 value. Inference requires $\ge 6$-week detraining episode (§11.6 result 3). \\
$S_\text{dec}$ & $0.07$ & S deconditioning threshold / §10 & Same. \\
$\mu_{B-}$ & $0.10$ & B deconditioning amplitude / §10 & Pinned to §10.4 Table 5 value. Inference requires same long detraining (§11.6 result 2). \\
$\mu_{S-}$ & $0.10$ & S deconditioning amplitude / §10 & Same. \\
\midrule
\multicolumn{4}{l}{\textit{Diffusion-side freezes (always frozen, both v4 and v5)}} \\
\midrule
$\sigma_B$ & $0.010$ & B Jacobi diffusion / §11.1 & Identifiability via diffusion is weak; standard practice. \\
$\sigma_S$ & $0.008$ & S Jacobi diffusion / §11.1 & Same. \\
$\sigma_F$ & $0.012$ & F CIR diffusion / §11.1 & Same. \\
$\sigma_A$ & $0.020$ & A CIR diffusion / §11.1 & Same. \\
$\sigma_K$ & $0.005$ & K CIR diffusion (shared) / §11.1 & Same. \\
$\phi$    & $0.0$    & circadian phase reference / §11.1 & Held at zero by convention. \\
\bottomrule
\end{tabular}
\end{center}

\subsection*{Initial state}

The 6D initial state is set from data or from \texttt{DEFAULT\_INIT}; it
has no explicit prior in the v5 \texttt{EstimationModel} (i.e.,
\texttt{INIT\_STATE\_PRIOR\_CONFIG = OrderedDict()}). Default values
correspond to a ``deconditioned but not unhealthy'' baseline:

\begin{center}\small
\begin{tabular}{lll}
\toprule
State & Default & Bounds \\
\midrule
$B$    & $0.05$  & $[0, 1]$ \\
$S$    & $0.10$  & $[0, 1]$ \\
$F$    & $0.30$  & $[0, \infty)$ \\
$A$    & $0.10$  & $[0, \infty)$ \\
$K_{FB}$ & $0.030$ & $[0, \infty)$ \\
$K_{FS}$ & $0.050$ & $[0, \infty)$ \\
\bottomrule
\end{tabular}
\end{center}

\subsection*{Reading guide for the smc2fc port}

After importing \texttt{models.fsa\_high\_res}, the port has access to:

\begin{itemize}
  \item \texttt{HIGH\_RES\_FSA\_V5\_ESTIMATION} — the ``EstimationModel`` to register with the SMC$^2$ engine. Built from the parameter prior config and frozen-params dict above.
  \item \texttt{HIGH\_RES\_FSA\_V5\_MODEL} — the ``SDEModel`` for forward simulation.
  \item \texttt{StepwisePlant} — the closed-loop plant for synthetic-feedback testing.
  \item \texttt{evaluate\_chance\_constrained\_cost} — the §9.6 v5 control cost (equations 23, 24).
\end{itemize}

The frozen values can be relaxed one at a time (move from
\texttt{frozen\_params} into \texttt{PARAM\_PRIOR\_CONFIG}) as the
calibration data set grows to support them.


\end{document}
